  % ======================================================================
  %  MCAV — Actuator-aware geometric flight control of a quadrotor
  %  Victor Gabriele De Matteis
  %
  %  Build:  (1) cd mcav_project && uv run python main.py
  %          (2) cd docs/report && latexmk -pdf report.tex
  %  Every number in this document is a macro from results/numbers.tex,
  %  generated by main.py; none is typed by hand.
  % ======================================================================

  \documentclass[10pt,a4paper]{scrartcl}

  % --------------------------------------------------
  % Encoding and language
  % --------------------------------------------------
  \usepackage[T1]{fontenc}
  \usepackage[utf8]{inputenc}
  \usepackage[english]{babel}

  % --------------------------------------------------
  % Page layout
  % --------------------------------------------------
  \usepackage[
    top=18mm,
    bottom=18mm,
    left=20mm,
    right=20mm,
    includefoot
  ]{geometry}

  % --------------------------------------------------
  % Mathematics and graphics
  % --------------------------------------------------
  \usepackage{graphicx}
  \usepackage{amsmath,amssymb}
  \usepackage{bm}
  \setlength{\jot}{0.25cm}
  \graphicspath{{../mcav_project/results/figures/}}
  \usepackage{tikz}
  \usepackage{float}
  \usetikzlibrary{positioning,arrows.meta}

  % --------------------------------------------------
  % Tables
  % --------------------------------------------------
  \usepackage{booktabs}
  \usepackage{tabularx}
  \usepackage{array}
  \usepackage{ragged2e}
  \newcolumntype{L}[1]{>{\RaggedRight\arraybackslash}p{#1}}

  % --------------------------------------------------
  % Typography
  % --------------------------------------------------
  \usepackage[scaled=0.95]{helvet}
  \usepackage[expansion=false]{microtype}
  \renewcommand{\familydefault}{\sfdefault}

  % --------------------------------------------------
  % Lists
  % --------------------------------------------------
  \usepackage{enumitem}
  \setlist[itemize]{leftmargin=*, topsep=2pt, parsep=0pt, partopsep=0pt, itemsep=3pt}
  \setlist[description]{leftmargin=0pt, topsep=2pt, parsep=0pt, itemsep=3pt}

  % --------------------------------------------------
  % Colour and hyperlinks
  % --------------------------------------------------
  \usepackage[dvipsnames]{xcolor}
  \usepackage[hidelinks]{hyperref}
  \definecolor{accent}{RGB}{0,129,88}
  \renewcommand{\emph}[1]{\textcolor{accent}{#1}}

  % --------------------------------------------------
  % Citations
  % --------------------------------------------------
  \usepackage[numbers,square,sort&compress]{natbib}

  % --------------------------------------------------
  % Header and footer
  % --------------------------------------------------
  \usepackage[automark,headsepline=false,footsepline=false]{scrlayer-scrpage}
  \clearpairofpagestyles
  \ifoot{Victor Gabriele De Matteis}
  \cfoot{\pagemark}
  \ofoot{MCAV --- Technical report}
  \setlength{\footskip}{20mm}
  \pagestyle{scrheadings}

  % --------------------------------------------------
  % Spacing and sections
  % --------------------------------------------------
  \setlength{\parindent}{0pt}
  \setlength{\parskip}{0.45em}
  \RedeclareSectionCommand[beforeskip=1.4ex plus 0.2ex minus 0.1ex, afterskip=0.6ex]{section}
  \RedeclareSectionCommand[beforeskip=1.0ex plus 0.2ex minus 0.1ex, afterskip=0.4ex]{subsection}
  \RedeclareSectionCommand[beforeskip=0.7ex plus 0.2ex minus 0.1ex, afterskip=0.3ex]{subsubsection}
  \setkomafont{section}{\large\bfseries\color{accent}}
  \setkomafont{subsection}{\normalsize\bfseries\color{accent}}
  \setkomafont{subsubsection}{\normalsize\bfseries}
  \newcommand{\RuleSection}[1]{%
    \section{#1}%
    \vspace{0.2em}{\color{accent!35}\hrule height 0.4pt}\vspace{0.3em}%
  }

  % --------------------------------------------------
  % Results generated by main.py
  % --------------------------------------------------
  \input{../mcav_project/results/numbers.tex}

  % ======================================================================
  \begin{document}
  % ======================================================================

  \begin{center}
    {\Huge\bfseries\color{accent} Actuator-Aware Geometric Flight Control of a Quadrotor}\\[0.25cm]
    {\small MCAV project --- Technical report --- October 2026}\\[0.15cm]
    {\small Victor Gabriele De Matteis}\\[0.2cm]
    {\color{accent!35}\hrule height 0.5pt}
  \end{center}

  \vspace{0.3em}

  \begin{abstract}
  \noindent
  This project delivers a flight controller for a 0.8\,kg quadrotor that takes off and tracks a
  circular trajectory in the MuJoCo physics engine. The controller works directly on rotation
  matrices (no Euler-angle singularities). Its gains are not hand-tuned: they are computed by
  convex optimisation (linear matrix inequalities) so that every motor stays inside its
  0--8\,N range. On the nominal mission the tracking error converges to $\RFinalErr$\,m and the
  motor forces stay between \RFMin\ and \RFMax\,N. With gains taken from the literature, the
  same mission would need between \RPaperFMin\ and \RPaperFMax\,N per motor. A Monte Carlo campaign over
  \RMcN\ random initial positions and attitudes (up to 175$^\circ$) converges in \RMcConverged\,\% of the cases.
  Every failure there comes from hitting the ground, which the analysed model does not contain.
  Stability claims are checked by a computer-algebra system instead of being asserted by hand,
  and every number in this report is regenerated by one command.
  \end{abstract}

  \begin{table}[h]
  \small
  \centering
  \renewcommand{\arraystretch}{1.1}
  \begin{tabularx}{\textwidth}{@{}L{6.2cm} X@{}}
  \toprule
  \textbf{Key result} & \textbf{Value} \\[.15cm]
  \midrule
  Steady-state tracking error & $\RFinalErr$\,m (settles below 5\,cm in \RSettleTime\,s) \\[.15cm]
  Motor forces, nominal mission & \RFMin--\RFMax\,N, inside the 0--8\,N motor range \\[.15cm]
  Same mission, gains from the literature & \RPaperFMin\ to \RPaperFMax\,N, not realisable \\[.15cm]
  Monte Carlo (\RMcN\ runs, attitude errors up to 175$^\circ$) & \RMcConverged\,\% converge; \RMcFeasible\,\% never leave the motor range \\[.15cm]
  Inside the certified attitude region (\RMcInside\ runs) & \RMcInsideConv\,\% converge, as the analysis predicts \\[.15cm]
  Symbolic identities verified by SymPy & \RNChecks \\[.15cm]
  Gain synthesis (6 per-axis LMI designs, open-source solver) & \RLmiTime\,s \\[.15cm]
  Full reproduction of every result & one command, \RRuntime\,s on a laptop \\[.15cm]
  \bottomrule
  \end{tabularx}
  \end{table}

  % ======================================================================
  \RuleSection{Problem and requirements}
  \label{sec:problem}

  A quadrotor of mass 0.8\,kg must take off from rest at $(0, 0, 0.5)$\,m and track the
  horizontal circle
  \[
    p_d(t) = [\,3\cos(0.5\,t),\ 3\sin(0.5\,t),\ 3\,]^\top\ \text{m},
  \]
  heading along the tangent.
  Each of its four motors delivers a thrust between 0 and 8\,N (hover needs 1.96\,N per motor).
  The physics is integrated by MuJoCo~\citep{todorov2012mujoco} at 500\,Hz, and the controller
  runs at the same rate.

  \begin{table}[h]
  \small
  \centering
  \renewcommand{\arraystretch}{1.1}
  \begin{tabularx}{\textwidth}{@{}L{0.7cm} L{6.3cm} X@{}}
  \toprule
  \textbf{ID} & \textbf{Requirement} & \textbf{How it is met} \\[.15cm]
  \midrule
  R1 & Zero steady-state tracking error & Exact feedforward of the reference acceleration and of the desired-frame motion (Sec.~\ref{sec:control}) \\[.15cm]
  R2 & Motor forces inside 0--8\,N on the nominal mission & Gains synthesised from the motor budget by LMIs (Sec.~\ref{sec:lmi}) \\[.15cm]
  R3 & Every stability claim backed by checkable evidence & Proofs verified by SymPy, certificates re-checked by tests (Sec.~\ref{sec:vv}) \\[.15cm]
  R4 & Robustness to the initial condition quantified & Monte Carlo over position and attitude errors (Sec.~\ref{sec:results}) \\[.15cm]
  R5 & Reproducible by a third party with one command & \texttt{uv run python main.py} (Sec.~\ref{sec:run}) \\[.15cm]
  \bottomrule
  \end{tabularx}
  \end{table}

  % ======================================================================
  \RuleSection{Architecture}
  \label{sec:arch}

  The controller is a two-level cascade (Fig.~\ref{fig:arch}). The outer loop turns the position
  error into a desired thrust vector. Its direction, together with the desired heading, defines
  the desired attitude $R_d$. The inner loop computes the torques that rotate the body onto
  $R_d$. A fixed allocation matrix then converts total thrust and torques into the four motor
  forces.

  \begin{figure}[h]
  \centering
  \resizebox{\textwidth}{!}{%
  \begin{tikzpicture}[
      node distance=0.5cm and 0.62cm,
      blk/.style={draw=accent, rounded corners=2pt, align=center, font=\footnotesize, minimum height=0.95cm, inner sep=3pt},
      arr/.style={-{Latex[length=1.6mm]}, thick, accent!80!black}]
    \node[blk] (ref) {Reference\\$p_d,\dot p_d,\ddot p_d,\psi_d$};
    \node[blk, right=of ref] (pos) {Position loop\\PD + feedforward};
    \node[blk, right=of pos] (att) {Desired attitude\\$R_d,\ \Omega_d,\ \dot\Omega_d$};
    \node[blk, right=of att] (in) {Attitude loop\\geometric, on $SO(3)$};
    \node[blk, right=of in] (alloc) {Allocation\\$f = A^{-1}[T,\tau]$};
    \node[blk, below=0.6cm of in, minimum width=4.2cm] (plant) {MuJoCo plant (0.8\,kg rigid body, 4 motors)};
    \draw[arr] (ref) -- (pos);
    \draw[arr] (pos) -- node[above=1pt, font=\scriptsize]{$F_{\mathrm{d}}$} (att);
    \draw[arr] (att) -- (in);
    \draw[arr] (in) -- node[above, font=\scriptsize]{$\tau$} (alloc);
    \draw[arr] (alloc) |- (plant);
    \draw[arr] (plant) -| node[pos=0.25, below, font=\scriptsize]{$p, v, R, \Omega$} (pos);
    \draw[arr] (plant) -- (in);
  \end{tikzpicture}}
  \caption{Control architecture. $F_{\mathrm{d}}$: desired thrust vector; its magnitude $T$ goes directly to the allocation.}
  \label{fig:arch}
  \end{figure}

  \textbf{Software.} The guiding rule was to write only the parts that are specific to this
  problem and to rely on maintained libraries for everything else (Table~\ref{tab:stack}).
  Vehicle parameters have a single source of truth: mass, inertia, motor positions, rotor
  torque coefficient and motor range are read from the MuJoCo model file. The allocation matrix
  in particular is MuJoCo's own actuator transmission, not a hand-written matrix.

  \begin{table}[h]
  \small
  \centering
  \renewcommand{\arraystretch}{1.1}
  \begin{tabularx}{\textwidth}{@{}L{4.2cm} L{3.4cm} X@{}}
  \toprule
  \textbf{Need} & \textbf{Library} & \textbf{Role} \\[.15cm]
  \midrule
  Rigid-body physics, allocation & MuJoCo & Simulation; mass, inertia and allocation read from the model \\[.15cm]
  Symbolic verification & SymPy & Derives and checks every identity of the stability analysis \\[.15cm]
  Gain synthesis & cvxpy + Clarabel~\citep{diamond2016cvxpy} & Semidefinite programs, open-source solver, no licence \\[.15cm]
  Rotations, ODE integration & SciPy & Attitude sampling, unicycle integration (RK45) \\[.15cm]
  Parallel Monte Carlo & joblib & One rollout per core \\[.15cm]
  Data, plots, tests & pandas, matplotlib, pytest & Logs, figures, 29 automated tests \\[.15cm]
  Environment & uv & Locked, reproducible dependency set \\[.15cm]
  \bottomrule
  \end{tabularx}
  \caption{Off-the-shelf components. The project-specific code (controller, gain design, analysis, experiments) is about 1\,100 lines of Python.}
  \label{tab:stack}
  \end{table}

  % ======================================================================
  \RuleSection{Control design}
  \label{sec:control}

  The design follows the geometric tracking controller of Lee et
  al.~\citep{lee2010geometric}, with the weighted attitude error of~\citep{lee2012exponential}
  so that every axis can get its own gain.

  \textbf{Position loop.} With $e_p = p - p_d$ and $e_v = \dot p - \dot p_d$,
  \begin{equation}
    F_{\mathrm{des}} = m\bigl(-K_v K e_p - K_v e_v + \ddot p_d + g e_3\bigr), \qquad T = \|F_{\mathrm{des}}\|.
  \end{equation}
  The feedforward term $\ddot p_d$ supplies the centripetal acceleration. Without it, the vehicle
  settles on a smaller, lagging circle, with a 17\,cm error in the previous version of this project.

  \textbf{Attitude loop.} $R_d$ has its third axis along $F_{\mathrm{des}}$ and its first axis as
  close as possible to the heading. With the attitude error $R_e = R_d^\top R$,
  \begin{equation}
    \tau = \underbrace{-e_R - K_\Omega e_\Omega}_{\text{feedback}}
           \;+\; \underbrace{\Omega\times J\Omega - J\bigl(\hat\Omega R_e^\top\Omega_d - R_e^\top\dot\Omega_d\bigr)}_{\text{cancellation}},
    \qquad e_R = \tfrac12\bigl(GR_e - R_e^\top G\bigr)^\vee .
  \end{equation}
  The cancellation terms remove the gyroscopic coupling and the motion of the desired frame.
  Without them the residual error does not vanish (\RPdRms\,mm RMS on the nominal mission) and
  the stability proof does not apply. The desired angular velocity $\Omega_d$ and its derivative
  come from the sequence of desired attitudes by finite differences. On the steady circle this is
  exact; during transients it is accurate to first order in the 2\,ms control period.

  % ======================================================================
  \RuleSection{Actuator-aware gain synthesis}
  \label{sec:lmi}

  Gains published for this controller were designed for a 4.3\,kg vehicle with an inertia about
  20 times larger. On this vehicle they demand \RPaperFMin\ to \RPaperFMax\,N per motor
  (Fig.~\ref{fig:forces}, left), far outside what the motors can deliver. Rather than re-tuning by
  trial and error, the gains are computed from the motor budget.

  \textbf{Budget.} Each motor force is $f_i = T/4 + \delta_i$, where $\delta_i$ is set by the
  torques. Reserving 1.5\,N of each motor's range for torques, and splitting it among roll,
  pitch and yaw, gives torque budgets of \RTauBudgetXY\,N\,m (roll, pitch) and \RTauBudgetZ\,N\,m
  (yaw: the rotor drag torque is small, so yaw authority is about ten times weaker). The remaining
  range bounds the correction acceleration to \RAccelBudgetXY\,m/s$^2$ horizontally and
  \RAccelBudgetZ\,m/s$^2$ vertically. If every loop respects its budget, every motor stays
  inside 0--8\,N.

  \textbf{Synthesis.} For each axis of each loop (a double integrator once the cancellation terms
  act) we solve the input-bounded state-feedback problem~\citep{boyd1994lmi}: find $W\succ 0$
  and $X$ with $K = XW^{-1}$ such that
  \begin{equation}
    AW + WA^\top + BX + X^\top B^\top \preceq -2\alpha W, \qquad
    \begin{bmatrix} 1 & x_0^\top\\ x_0 & W\end{bmatrix}\succeq 0, \qquad
    \begin{bmatrix} \bar u^2 & X\\ X^\top & W\end{bmatrix}\succeq 0,
  \end{equation}
  and maximise the decay rate $\alpha$ by bisection. Every trajectory of the linear model that
  starts in the ellipsoid $\{x^\top W^{-1}x\le 1\}$ stays in it, decays at rate $\alpha$, and
  never asks for more than $\bar u$. The design points $x_0$ are the mission start plus a margin.
  The result is a position loop with rates of \RAlphaPosMin--\RAlphaPosMax\,s$^{-1}$ and an
  attitude loop at \RAlphaRollPitch\,s$^{-1}$ (roll, pitch) and \RAlphaYaw\,s$^{-1}$ (yaw). The
  inner loop is therefore faster than the outer loop, as the cascade needs. The nominal mission
  then stays between \RFMin\ and \RFMax\,N per motor (Fig.~\ref{fig:forces}, right), at the cost
  of a slower transient: \RSettleTime\,s instead of \RPaperSettle\,s to reach 5\,cm.

  \begin{figure}[h]
  \centering
  \includegraphics[width=\textwidth]{motor_forces}
  \caption{Motor forces during take-off. Left: gains from the literature. Right: LMI gains. Shaded: motor range.}
  \label{fig:forces}
  \end{figure}

  % ======================================================================
  \RuleSection{Verification and validation}
  \label{sec:vv}

  The project separates what is proved, what is checked numerically, what is observed in
  simulation and what is assumed (Table~\ref{tab:evidence}).

  \begin{table}[h]
  \small
  \centering
  \renewcommand{\arraystretch}{1.1}
  \begin{tabularx}{\textwidth}{@{}L{1.7cm} X L{4.6cm}@{}}
  \toprule
  \textbf{Status} & \textbf{Statement} & \textbf{Evidence} \\[.15cm]
  \midrule
  Proved & The attitude error converges from every state with $V_2 < \RLevelBound$, i.e.\ roll/pitch errors up to about \RCertRollDeg$^\circ$ and any yaw error below 180$^\circ$ at zero initial rate & Lyapunov function and LaSalle; \RNChecks\ identities verified by SymPy \\[.15cm]
  Proved & The full closed loop is locally exponentially stable & Cascade structure and linearisation~\citep{khalil2002nonlinear} \\[.15cm]
  Proved & On the reduced linear models, the LMI gains respect the motor budget inside the certified ellipsoids & LMI certificates, re-checked by tests independently of the solver \\[.15cm]
  Observed & Convergence and motor feasibility from large initial errors & MuJoCo simulation, Monte Carlo \\[.15cm]
  Assumed & Motors can produce any force (no saturation in the analysis); $\Omega_d,\dot\Omega_d$ exact (implemented by finite differences) & Feasibility is reported when the range would be exceeded \\[.15cm]
  \bottomrule
  \end{tabularx}
  \caption{Evidence behind each claim.}
  \label{tab:evidence}
  \end{table}

  \textbf{What is not claimed.} The closed loop is not globally asymptotically stable. No
  continuous feedback on the rotation group can be~\citep{bhat2000topological}: a 180$^\circ$
  attitude error is an equilibrium of any such controller. The analysis states its region
  explicitly instead.

  \textbf{Symbolic verification.} The stability argument rests on a few exact identities: the
  attitude error function decreases at the rate $\dot\Psi = e_R\cdot e_\Omega$; the cancellation
  terms make the attitude error dynamics independent of the trajectory; and the only attitudes
  with $e_R = 0$ are the identity and 180$^\circ$ rotations. SymPy derives each identity from the
  equations of motion and the control law, and checks it for a generic matrix, so no small-angle
  approximation is involved. It also produces a concrete state at which the feedback-only law
  makes the Lyapunov function increase. That state shows why the cancellation terms are needed.

  \textbf{Automated tests.} 29 \texttt{pytest} tests cover the symbolic identities, the LMI
  certificates (decay rate, design points inside the ellipsoid, input bound over the ellipsoid),
  and the closed loop. Among them: the mission converges, all motor forces stay in range, the
  Lyapunov function never increases along the simulated trajectory, the finite-difference
  $\Omega_d$ matches its exact steady-state value, and the feedback-only law leaves a residual error.

  % ======================================================================
  \RuleSection{Results}
  \label{sec:results}

  \textbf{Nominal mission.} The vehicle converges to the circle and the error decreases to
  machine precision (Fig.~\ref{fig:mission}). The feedback-only law, without the cancellation
  terms, stalls at the millimetre level. The attitude Lyapunov function never increases along
  the simulated trajectory (Fig.~\ref{fig:lyap}), as the analysis predicts.

  \begin{figure}[h]
  \centering
  \includegraphics[width=\textwidth]{mission_tracking}
  \caption{Nominal mission with the LMI gains. Left: top view. Right: position error, with and without the cancellation terms; dotted line at 5\,cm.}
  \label{fig:mission}
  \end{figure}

  \textbf{Robustness to the initial condition.} \RMcN\ runs start from random position offsets
  (up to 2\,m) and random attitude errors (up to 175$^\circ$ about a random axis), at rest, and
  fly the same mission with the same gains (Fig.~\ref{fig:mc}). \RMcConverged\,\% converge.
  The other \RMcGround\,\% hit the ground during the recovery manoeuvre; the floor exists in the
  simulation but not in the analysed model. All \RMcInside\ runs that start inside the certified
  attitude region converge, as the analysis predicts. \RMcFeasible\,\% of all runs never need more than the motors can
  deliver, even though the LMI certificate covered only the nominal take-off envelope.

  \begin{figure}[h]
  \centering
  \begin{minipage}[t]{0.55\textwidth}
    \centering
    \includegraphics[width=\textwidth]{monte_carlo}
    \caption{Monte Carlo: outcome for each initial condition.}
    \label{fig:mc}
  \end{minipage}\hfill
  \begin{minipage}[t]{0.42\textwidth}
    \centering
    \includegraphics[width=\textwidth]{lyapunov_attitude}
    \caption{Attitude Lyapunov function on the nominal mission.}
    \label{fig:lyap}
  \end{minipage}
  \end{figure}

  % ======================================================================
  \RuleSection{Engineering review: defects found and fixed}
  \label{sec:review}

  A systematic review of an earlier version of the project found the defects below. Each one is
  now covered by a fix and, where possible, by an automated test.

  \begin{table}[H]
  \small
  \centering
  \renewcommand{\arraystretch}{1.1}
  \begin{tabularx}{\textwidth}{@{}L{5.0cm} L{4.2cm} X@{}}
  \toprule
  \textbf{Defect} & \textbf{Measured impact} & \textbf{Fix} \\[.15cm]
  \midrule
  Reference acceleration not fed forward & 17\,cm steady-state error & Exact $\ddot p_d$ from the reference \\[.15cm]
  Moment arm set to the arm length $L$ instead of $L/\sqrt2$ & Roll/pitch torques realised at 71\,\% & Allocation read from the MuJoCo model \\[.15cm]
  Simulator silently clipped motor commands, while the analysis assumed no limits & Simulated and analysed systems differed & Limits stored in the model, enforced as a design constraint, reported in validation \\[.15cm]
  Desired angular velocity approximated by the yaw rate & 3\,cm residual error even with the acceleration fixed & $\Omega_d$, $\dot\Omega_d$ from $R_d$ \\[.15cm]
  Gains copied from a 4.3\,kg vehicle & Peaks above 150\,N per motor & LMI synthesis from the motor budget \\[.15cm]
  ``Globally asymptotically stable'' claimed with a hand-written derivative & Unsupported, and impossible on $SO(3)$ & Verified local results with explicit regions \\[.15cm]
  About 900 lines of custom simulation, logging and plotting code & Maintenance cost & Replaced by MuJoCo, SciPy, joblib, pandas \\[.15cm]
  \bottomrule
  \end{tabularx}
  \end{table}

  % ======================================================================
  \RuleSection{Limitations and next steps}
  \label{sec:next}

  \begin{itemize}
    \item \textbf{Actuator saturation.} Motor limits are respected by design on the nominal
      mission, but are not enforced online. Next step: a control barrier function or a model
      predictive controller acting on the same allocation.
    \item \textbf{Unmodelled effects.} Motor dynamics, delays, aerodynamic drag, ground effect
      and sensor noise are not simulated. The state is assumed to be measured exactly. Next
      step: an estimator (EKF on IMU data) and a robustness study.
    \item \textbf{Path to hardware.} The controller is a pure function of state and reference
      at a fixed rate. It can be ported to a PX4 or ROS~2 software-in-the-loop setup before any
      flight test.
  \end{itemize}

  % ======================================================================
  \RuleSection{Reproducing the results}
  \label{sec:run}

  The only prerequisite is \texttt{uv}, which installs Python and the locked dependency set.
  \begin{description}
    \item[] \texttt{cd mcav\_project \&\& uv run python main.py} --- runs the symbolic checks,
      the LMI design, the simulations and the Monte Carlo, and writes the figures and
      \texttt{numbers.tex} used by this report (\RRuntime\,s).
    \item[] \texttt{uv run pytest} --- runs the automated tests.
    \item[] \texttt{uv run jupyter lab} --- opens \texttt{notebooks/final.ipynb}, a guided
      walkthrough of the same steps.
  \end{description}

  \bibliographystyle{abbrvnat}
  \bibliography{references}

  \end{document}
